\section{Derivation of servo coefficients}
\label{sec:appendix_servo}

This appendix derives the servo coefficients used in the combined servo mechanism. The torsional servo coefficient $S_T$ is independent of the other conditions and has been defined in \cref{eq:irot_scs}. Below, the radial and axial servo coefficients are derived by considering the volume-coupling effect.

From the constant-volume constraint (\cref{eq:const_volume_constraint}), the outer radius rate is determined by:
\begin{equation}
\frac{dR}{dt} = \frac{r}{R}\frac{dr}{dt} - \frac{R^2-r^2}{2RH}\frac{dH}{dt}.
\label{eq:dRdt_general}
\end{equation}

Movement of both the inner wall ($dr$) and the top/bottom plates ($dH$) induces a change in the outer radius through \cref{eq:dRdt_general}, which in turn affects both $p_{dif,r}^{\prime}$ and $p_{dif,z}^{\prime}$. Although the vertical pressure $p_z^{\prime}$ depends only on the top/bottom contacts, the axial pressure difference $p_{dif,z}^{\prime}=p_z^{\prime}-\sigma_r^{\prime}$ is coupled to radial wall movements through $\sigma_r^{\prime}$.

Let $K_{r,i}=\sum_{c\in\mathrm{inner}} k_n$ and $K_{r,o}=\sum_{c\in\mathrm{outer}} k_n$ denote the aggregate normal contact stiffness at the inner and outer cylinders, respectively, and $K_z=\frac{1}{2}\sum_{c\in\mathrm{top/bottom}} k_n$ the effective axial contact stiffness, where the summations run over all particle--wall contacts on the respective rigid boundaries. The factor $\frac{1}{2}$ arises because the top and bottom plates move symmetrically (each by $dH/2$) to achieve the total height change $dH$. The corresponding wall areas are $A_i=2\pi rH$, $A_o=2\pi RH$, and $A_z=2\pi(R^2-r^2)$ (see Table~\ref{tab:hca_equations}).

The sign convention adopts outward-positive for all kinematic increments: $dr>0$ and $dR>0$ denote outward (radius-increasing) wall movements, and $dH>0$ denotes an increase in specimen height. The individual pressure increments are:
\begin{align}
dp_i^{\prime} &= \frac{K_{r,i}}{A_i}\,dr,\qquad
dp_o^{\prime} = -\frac{K_{r,o}}{A_o}\,dR, \label{eq:dp_individual_r}\\[4pt]
dp_z^{\prime} &= -\frac{K_z}{A_z}\,dH, \label{eq:dp_individual_z}
\end{align}
where the signs follow from the geometry: an outward movement of the inner wall ($dr>0$) expands the inner cylinder into the specimen, compressing particles against it and increasing $p_i^{\prime}$; an outward movement of the outer wall ($dR>0$) moves it away from the specimen, unloading particles and decreasing $p_o^{\prime}$; and an increase in specimen height ($dH>0$) moves the top and bottom plates apart, unloading axial contacts and decreasing $p_z^{\prime}$.

The radial pressure difference $p_{dif,r}^{\prime}=p_o^{\prime}-p_i^{\prime}$ and the average radial effective stress $\sigma_r^{\prime}=({p_o^{\prime}R+p_i^{\prime}r})/({R+r})$ (Table~\ref{tab:hca_equations}) then give:
\begin{align}
dp_{dif,r}^{\prime}
  &= dp_o^{\prime}-dp_i^{\prime}
   = -\frac{K_{r,i}}{A_i}\,dr - \frac{K_{r,o}}{A_o}\,dR, \label{eq:dp_r_three}\\[6pt]
d\sigma_r^{\prime}
  &= \frac{r}{R+r}\,dp_i^{\prime} + \frac{R}{R+r}\,dp_o^{\prime}
   = \frac{K_{r,i}\,dr - K_{r,o}\,dR}{A_i+A_o}, \label{eq:dsigma_r_three}\\[6pt]
dp_{dif,z}^{\prime}
  &= dp_z^{\prime}-d\sigma_r^{\prime}
   = -\frac{K_z}{A_z}\,dH - \frac{K_{r,i}\,dr - K_{r,o}\,dR}{A_i+A_o}. \label{eq:dp_z_three}
\end{align}
Note that both terms in \cref{eq:dp_r_three} carry negative signs: any outward wall movement---whether inner or outer---reduces $p_{dif,r}^{\prime}$. By contrast, $d\sigma_r^{\prime}$ involves the weighted average from Table~\ref{tab:hca_equations}, in which the geometric factors $r/(R+r)$ and $R/(R+r)$ weight the inner and outer contributions differently.

Substituting \cref{eq:dRdt_general} into \cref{eq:dp_r_three,eq:dp_z_three} to eliminate $dR$ yields the pressure responses in terms of only $dr$ and $dH$:
\begin{align}
dp_{dif,r}^{\prime}
  &= -\!\left(\frac{K_{r,i}}{A_i}+\frac{r}{R}\,\frac{K_{r,o}}{A_o}\right)dr
     +\frac{R^2-r^2}{2RH}\,\frac{K_{r,o}}{A_o}\,dH, \label{eq:dp_r_two}\\[6pt]
dp_{dif,z}^{\prime}
  &= -\frac{K_{r,i}-\frac{r}{R}\,K_{r,o}}{A_i+A_o}\,dr
     -\!\left(\frac{K_z}{A_z}+\frac{(R^2-r^2)\,K_{r,o}}{2RH\,(A_i+A_o)}\right)dH. \label{eq:dp_z_two}
\end{align}
Collecting the coefficients of $dr$ and $dH$ into a matrix:
\begin{equation}
\begin{pmatrix} dp_{dif,r}^{\prime} \\[4pt] dp_{dif,z}^{\prime} \end{pmatrix}
=
\underbrace{\begin{pmatrix} \hat{K}_{11} & \hat{K}_{12} \\[8pt] \hat{K}_{21} & \hat{K}_{22} \end{pmatrix}}_{\hat{\mathbf{K}}}
\begin{pmatrix} dr \\[4pt] dH \end{pmatrix},
\label{eq:stiffness_matrix}
\end{equation}
where
\begin{alignat}{2}
\hat{K}_{11} &= -\frac{K_{r,i}}{A_i}-\frac{r}{R}\,\frac{K_{r,o}}{A_o}, &\qquad
\hat{K}_{12} &= \frac{R^2-r^2}{2RH}\,\frac{K_{r,o}}{A_o}, \label{eq:Khat_row1}\\[4pt]
\hat{K}_{21} &= -\frac{K_{r,i}-\frac{r}{R}\,K_{r,o}}{A_i+A_o}, &\qquad
\hat{K}_{22} &= -\frac{K_z}{A_z}-\frac{(R^2-r^2)\,K_{r,o}}{2RH\,(A_i+A_o)}. \label{eq:Khat_row2}
\end{alignat}
All entries carry dimensions of stress\,/\,length (Pa/m), reflecting the displacement-to-pressure mapping through the wall areas $A_i$, $A_o$, and $A_z$. Three entries are negative ($\hat{K}_{11}$, $\hat{K}_{21}$, $\hat{K}_{22}<0$), consistent with the outward-positive convention: outward wall movements generally reduce the confining pressures. The exception is $\hat{K}_{12}>0$, because an increase in specimen height forces the outer wall inward through the volume constraint (\cref{eq:dRdt_general}), which raises $p_o^{\prime}$ and thereby increases $p_{dif,r}^{\prime}$. Unlike the radial pressure difference, which receives equal-sign contributions from both cylinders (\cref{eq:dp_r_three}), the average radial stress $\sigma_r^{\prime}$ is a weighted mean of $p_o^{\prime}$ and $p_i^{\prime}$ (Table~\ref{tab:hca_equations}), so $\hat{K}_{21}\neq\hat{K}_{11}$ in general.

The determinant of $\hat{\mathbf{K}}$ is:
\begin{equation}
\det(\hat{\mathbf{K}}) = \hat{K}_{11}\,\hat{K}_{22} - \hat{K}_{12}\,\hat{K}_{21},
\label{eq:det_K}
\end{equation}
and the inverse is:
\begin{equation}
\hat{\mathbf{K}}^{-1}
=
\frac{1}{\det(\hat{\mathbf{K}})}
\begin{pmatrix} \hat{K}_{22} & -\hat{K}_{12} \\[6pt] -\hat{K}_{21} & \hat{K}_{11} \end{pmatrix}.
\label{eq:stiffness_matrix_inv}
\end{equation}

The servo must apply wall movements that eliminate the current pressure errors $e_r$ and $e_z$, i.e., produce $dp_{dif,r}^{\prime}=-e_r$ and $dp_{dif,z}^{\prime}=-e_z$. Setting $(dp_{dif,r}^{\prime},\,dp_{dif,z}^{\prime})=(-e_r,\,-e_z)$ in \cref{eq:stiffness_matrix} and solving for the displacement rates gives:
\begin{align}
\frac{dr}{dt} &= \frac{1}{\det(\hat{\mathbf{K}})}\!\left(-\hat{K}_{22}\,\frac{e_r}{\Delta t} + \hat{K}_{12}\,\frac{e_z}{\Delta t}\right), \label{eq:dr_general}\\[6pt]
\frac{dH}{dt} &= \frac{1}{\det(\hat{\mathbf{K}})}\!\left(\hat{K}_{21}\,\frac{e_r}{\Delta t} - \hat{K}_{11}\,\frac{e_z}{\Delta t}\right), \label{eq:dH_general}
\end{align}
where $e_r=p_{dif,r}^{\prime}-p_{dif,r}^{tar}$ and $e_z=p_{dif,z}^{\prime}-p_{dif,z}^{tar}$ are the pressure errors defined in \cref{eq:cond1_radial,eq:cond2_axial}.

Under the constant-height condition ($dH/dt=0$), the axial pressure difference $p_{dif,z}^{\prime}$ is not actively controlled, and only the radial condition remains. Setting $dH=0$ in \cref{eq:stiffness_matrix}, the radial servo coefficient reduces to:
\begin{equation}
S_{rr} = -\frac{1}{\hat{K}_{11}} = \frac{1}{\dfrac{K_{r,i}}{A_i}+\dfrac{r}{R}\,\dfrac{K_{r,o}}{A_o}}.
\label{eq:Srr_constantH}
\end{equation}

For the general case ($dH/dt\neq 0$), the fully coupled servo equations (\cref{eq:dr_general,eq:dH_general}) are applied at each timestep.
